\documentclass[a4paper,12pt]{jsarticle}
\usepackage[utf8]{inputenc}
\usepackage{amsmath,amssymb}
\usepackage{enumitem}
\usepackage{mdframed} 
\usepackage{bm}
\usepackage{physics}
\usepackage{float}
\usepackage{mathdots}
\usepackage[dvipdfm]{graphicx}
\newcommand{\p}{\partial}
\newcommand{\poi}[1]{\{ #1 \}}
%\newcommand{\log}{\mathrm{log}}
\newcommand{\ve}{\varepsilon}
\newcommand{\tdv}[2]{\frac{d #1}{d #2}}
\renewcommand{\pdv}[2]{\frac{\p #1}{\p #2}}
\renewcommand{\abs}[1]{| #1 |}
\newcommand{\Biglr}[1]{\Bigl[ #1 \Bigr]}
\newcommand{\Biglrm}[1]{\Bigl( #1 \Bigr)}
\renewcommand{\fdv}[2]{\frac{d #1}{d #2}}
\renewcommand{\order}[1]{\mathcal{O}[ #1 ]}
\newcommand{\m}[1]{\mathcal{#1}}
\newcommand{\h}[1]{\hat{#1}}
%\newcommand{\dd}{\mathrm{d}}


\begin{document}

\section{BosonとFermionのCoherent state}
このノートでは経路積分形式で使うためのcohenret状態とその性質についてまとめる。coherent状態を使う理由としては、後に見てわかる通り、場の演算子を微分・積分可能な数に置き換えるためである。

統計力学において、分配関数による物理量の計算を学んだことを覚えているだろうか？分配関数は一度積分を計算しておけば、あとは任意の物理量が微分(変分)によって計算が出来ることが利点であった。古典の場合は積分すれば良いのだが、量子の場合は分配関数の中身が演算子になっており、量子の演算子は一般に飛び飛びの値を取るため、そのままでは微分(変分)ができない。そこでcoherent状態を導入することで、演算子を微分・積分可能な数に置き換えることができるのである。(量子系を「古典系っぽく」扱うための方法だと思ってもらって良い。)

\subsection{ボゾンのcoherent状態}
場の生成・消滅演算子$\poi{a^{(\dagger)}_i}_{i=1}^N$が定義されている状態空間を考える。($[a_i, a_j^\dagger] = \delta_{ij}$が満たされているとする。)
この時コヒーレント状態はベクトル変数$\vec{\phi} = [\phi_1, \phi_2, \dots \phi_N]$を用いて
\begin{eqnarray}
  \ket{\vec{\phi}} = \exp\Biglr{\sum_i \phi_i a_i^\dagger}\ket{0} = \prod_i \exp\Biglr{\phi_i a_i^\dagger}\ket{0}
\end{eqnarray}
と定義される。これは以下を満たす。
\begin{eqnarray}
  a_i \ket{\vec{\phi}} = \phi_i \ket{\vec{\phi}}, \quad \bra{\vec{\phi}} a_i^\dagger = \bra{\vec{\phi}} \phi_i^* \label{coherent_eigen}
\end{eqnarray}
よってこの状態の上では演算子は$c$数(複素数)として振る舞う。

[練習問題]式(\ref{coherent_eigen})を示せ。



\subsection{Floquetの擬エネルギーとFloquet状態}
\begin{enumerate}
  \item (☆)$\h{F}(t_0)$の固有値$\epsilon_{\alpha}$の固有状態を$\ket{\psi_\alpha(t_0)}$とすると、
\begin{eqnarray}
  \ket{\psi_\alpha(t_0+T)} = \hat{U}(t_0+T,t_0)\ket{\psi_\alpha(t_0)} = e^{-i\hat{F} T}\ket{\psi_\alpha(t_0)} = e^{-i\epsilon_\alpha T}\ket{\psi_\alpha(t_0)}
\end{eqnarray}
を満たす、これが任意の$t$についても同様に$$\ket{\psi_\alpha(t+T)} = e^{-i\epsilon_\alpha T}\ket{\psi_\alpha(t)}$$を満たす事を示せ。よって、$\ket{\psi_\alpha(t)}$は時間$T$の離散並進操作$e^{-i\h{F}T}$の固有状態である事がわかる。(教科書の式(5.26)の時間版)
  \item (☆)さらに、上の結果を用いて($\ket{\psi_\alpha(t)}$を用いて)、Bloch状態(教科書の式(5.36))の時間版$\ket{v_\alpha(t)}$を導出せよ。また$\ket{v_\alpha(t+T)} = \ket{v_\alpha(t)}$を確認せよ。これがFloquet状態と呼ばれるものである。
  \item (☆☆)ちなみに$\h{F}$の表現は様々な形に変換する事が出来、用途によって使い分ける。適当な正規直交基底$\{\ket{\phi_\beta}\}$を用いて
\begin{eqnarray}
  \h{U}(t, t_0) &=& \sum_\alpha \ket{\psi_\alpha(t)}\bra{\psi_\alpha(t_0)} = \sum_\alpha \ket{v_\alpha(t)}e^{-i\epsilon_\alpha (t-t_0)}\bra{v_\alpha(t_0)}\\
  &=& \sum_{\alpha,\beta,\gamma} \ket{v_\alpha(t)}\delta_{\alpha\beta}e^{-i\epsilon_\beta (t-t_0)}\delta_{\beta\gamma}\bra{v_\gamma(t_0)}\\
  &=& \Biglrm{\sum_{\alpha} \ket{v_\alpha(t)}\bra{\phi_\alpha}} \Biglrm{\sum_\beta \ket{\phi_\beta}e^{-i\epsilon_\beta (t-t_0)}\bra{\phi_\beta}}\Biglrm{\sum_{\gamma}\ket{\phi_\gamma}\bra{v_\gamma(t_0)}}\\
  &=& \h{V}(t) e^{-i\h{H}_\mathrm{eff} (t-t_0)} \h{V}^\dagger(t_0)
\end{eqnarray}
と書き直せる。$\h{V}(t)$がユニタリ演算子である事を示せ。さらに
\begin{eqnarray}
  \h{H}_\mathrm{eff} = \h{V}^\dagger(t) \h{H}(t) \h{V}(t) - i \h{V}^\dagger(t) \pdv{\h{V}}{t}
\end{eqnarray}
と書ける事を示せ。これは任意の時間周期Hamiltonianに対して、それを時間依存しないHamiltonian $\h{H}_\mathrm{eff}$にMap出来るユニタリ演算子$\h{V}(t)$を見つける事が出来る事を意味する。(時間依存しないHamiltonianによって時間発展しているように見える観測系が存在する)
\end{enumerate}
\subsection{Sambe空間と高周波展開}
\begin{enumerate}
  \item (☆☆☆)今Floquet状態$\ket{v_\alpha(t)}$とHamiltonian$\h{H}(t)$は時間周期的であるから、フーリエ級数展開して
\begin{eqnarray}
  \ket{v_\alpha(t)} = \sum_{n=-\infty}^\infty e^{-in\Omega t}\ket{v_\alpha^n}, \quad \h{H}(t) = \sum_{m=-\infty}^\infty e^{-im\Omega t}\h{H}_m
\end{eqnarray}
と書ける。ただし$\Omega = 2\pi/T$である。これを用いて以下の式が成り立つ事を示せ。
\begin{eqnarray}
  \sum_{m=-\infty}^\infty \Biglrm{\h{H}_{n-m} - \delta_{nm}m\Omega}\ket{v_\alpha^{m}} = \epsilon_\alpha \ket{v_\alpha^n} \label{eff_Floquet}
\end{eqnarray}

\item (☆☆)さらに$\ket{\ket{v_\alpha}} = [\dots, \ket{v_\alpha^1}, \ket{v_\alpha^0}, \ket{v_\alpha^{-1}}, \dots]^T$で基底を貼られるSambe空間において、式(\ref{eff_Floquet})が、
\begin{eqnarray}
  \begin{pmatrix}
    \ddots & \vdots & \vdots & \vdots & \iddots\\
    \cdots & \h{H}_0 + \Omega & \h{H}_{-1} & \h{H}_{-2} & \cdots\\
    \cdots & \h{H}_1 & \h{H}_0 & \h{H}_{-1} & \cdots\\
    \cdots & \h{H}_2 & \h{H}_1 & \h{H}_0 -\Omega & \cdots\\
    \iddots & \vdots & \vdots & \vdots & \ddots
  \end{pmatrix}\begin{pmatrix}
    \vdots\\
    \ket{v_\alpha^1}\\
    \ket{v_\alpha^0}\\
    \ket{v_\alpha^{-1}}\\
    \vdots
  \end{pmatrix} = 
  e^{-i\epsilon_\alpha t}\begin{pmatrix}
    \vdots\\
    \ket{v_\alpha^1}\\
    \ket{v_\alpha^0}\\
    \ket{v_\alpha^{-1}}\\
    \vdots
\end{pmatrix}\label{Sambe}
\end{eqnarray}
と書ける事を示せ。
\item (☆☆) しかし、式(\ref{Sambe})は無限次元の行列であり、実際にこれを解くのは不可能である。そこで、$\Omega$が$H_n$に比べて十分大きい時に使える有効的なHamiltonianを導出しよう。$\h{V}_d(t) = \h{V}_0 \cos(\Omega t)$である時、$\h{H}_0, \h{H}_1, \h{H}_{-1}$を導け。また教科書2.6章で勉強した有効ハミルトニアンの方法を用いて、$\ket{v_\alpha^0}$の空間に射影した有効ハミルトニアンが
\begin{eqnarray}
  \h{H}_\mathrm{eff} = \h{H}_0 + \sum_n\frac{1}{2n\Omega}[\h{H}_{-n}, \h{H}_{n}] + \order{\frac{1}{\Omega^2}}
\end{eqnarray}
と書ける事を示せ。
\end{enumerate}
\end{document}